\documentclass[11pt,a4paper]{article}
\usepackage[T1]{fontenc}
\usepackage[utf8]{inputenc}
\usepackage{amsmath,amssymb,amsthm}
\usepackage[margin=1in]{geometry}
\usepackage[colorlinks=true,linkcolor=blue,urlcolor=blue]{hyperref}
\theoremstyle{plain}
\newtheorem{theorem}{Theorem}
\newtheorem{lemma}[theorem]{Lemma}
\newtheorem{proposition}[theorem]{Proposition}
\newtheorem{corollary}[theorem]{Corollary}
\theoremstyle{definition}
\newtheorem{remark}[theorem]{Remark}

\title{The empirical recurrence for OEIS A232147, proved by transfer matrix}
\author{Adrian Perez Fontelles\\ \small Independent researcher}
\date{3 September 2026}
\begin{document}
\maketitle

\begin{abstract}
OEIS A232147 counts arrays of width $4$ over $\{0,1,2\}$ in which every cell carrying one stated
value must have a neighbour carrying another. The entry carries an empirical recurrence of
order $32$ contributed by R.~H.~Hardin. It is true, and it is decidable rather than
empirical: a cell's neighbours lie in three consecutive rows, so the count is a walk count on
$S=6643$ states.
\end{abstract}

\noindent\small 2020 Mathematics Subject Classification. 05A15, 05C15, 15A18.\normalsize

\section{The conjecture}

OEIS A232147 is ``Number of n X 4 0..2 arrays with every 0 next to a 1 and every 1 next to a 2 horizontally, diagonally or antidiagonally.''. It has offset $1$ and begins
\[
15, 1299, 67725, 3409361, 177194147, 9181257593, 475203378037, 24603676419865,\ \dots
\]
The entry states, as a conjecture contributed by its author and never marked settled:
\begin{quote}
Empirical: a(n) = 35*a(n-1) +649*a(n-2) +10479*a(n-3) +50649*a(n-4) -148850*a(n-5) -4594482*a(n-6) -29446017*a(n-7) -109126913*a(n-8) -189857827*a(n-9) +786606763*a(n-10) +4637371935*a(n-11) +5864860064*a(n-12) -12212688755*a(n-13) -42763674878*a(n-14) -28294485410*a(n-15) +57477014048*a(n-16) +126547355711*a(n-17) +64552653143*a(n-18) -60155915760*a(n-19) -96981273076*a(n-20) -64212382080*a(n-21) -123120194790*a(n-22) -55873185285*a(n-23) -22972924737*a(n-24) +79085877984*a(n-25) +44270889165*a(n-26) +95126815696*a(n-27) +10764163372*a(n-28) +18110185584*a(n-29) -9929666880*a(n-30) -6320643840*a(n-31) +1858775040*a(n-32).
\end{quote}
As of the ``Last modified'' line on the live entry (May 11 2023 15:06:35, revision 6) this is still
recorded as empirical, and nothing on the entry records it as proved.

\section{What is being counted}

The array has $4$ cells across, with entries in $\{0,1,2\}$. The NEIGHBOURS of a cell are the
cells reached from it by the offsets
\[
(-1,-1), (-1,1), (0,-1), (0,1), (1,-1), (1,1),
\]
those lying inside the array. The entry asks that
\begin{itemize}
\item every cell carrying $0$ has at least one neighbour carrying $1$;
\item every cell carrying $1$ has at least one neighbour carrying $2$;

\end{itemize}
A cell whose value is not named on the left of any clause is unconstrained.

\section{Three rows decide one}

\begin{lemma}\label{lem:win}
Every offset above changes the row index by $-1$, $0$ or $+1$, so all the neighbours of a cell
of row $i$ lie in rows $i-1$, $i$ and $i+1$, and the condition for the whole of row $i$ is
decided by those three rows.
\end{lemma}

\begin{proof}
Immediate from the offsets listed, together with the rule that a neighbour outside the array is
not a neighbour.
\end{proof}

So the state is the pair (row above, current row), the row above allowed to be the symbol
$\bot$, and a step appends a row and settles the condition for the middle row of the window it
completes. The condition for the LAST row is settled by the end vector, with nothing below it:
the end vector is therefore not the all-ones vector, and a row whose cells find what they need
only because a further row is written below them does not finish an array. With $M$ the
adjacency matrix,
\[
a(n)\;=\;\iota^{\!\top}M^{\,j}\tau ,\qquad S=6643.
\]
The walk index $j$ and the entry's own $n$ differ by a constant, read off by matching the
published terms rather than assumed. In particular $a$ is $C$-finite.

\section{The proof}

\begin{lemma}\label{lem:crit}
Let $q(t)=t^{32}-\sum_i c_it^{32-i}$ be the characteristic polynomial of the
conjectured recurrence. The residual $a(n)-\sum_ic_ia(n-i)$ equals
$\iota^{\!\top}M^{\,j}q(M)\tau$ for the corresponding walk index $j$, so the recurrence
holds from some point on if and only if $\iota^{\!\top}M^{j}q(M)\tau=0$ for all large $j$,
and it is enough to examine $j<S$.
\end{lemma}

\begin{proof}
Factor $M^{j}$ out of $M^{j+32}-\sum_ic_iM^{j+32-i}$. For the bound, $M^{S}$
is an integer combination of $I,M,\dots,M^{S-1}$ by Cayley--Hamilton, so the numbers
$u_j=\iota^{\!\top}M^{j}q(M)\tau$ obey the monic recurrence given by the characteristic
polynomial of $M$; $S$ consecutive zeros force every later one, and the last nonzero $u_j$
pins the threshold exactly.
\end{proof}

\begin{theorem}
\[
a(n)\;=\;35\,a(n-1) + 649\,a(n-2) + 10479\,a(n-3) + 50649\,a(n-4) - 148850\,a(n-5) - 4594482\,a(n-6) - 29446017\,a(n-7) - 109126913\,a(n-8) - 189857827\,a(n-9) + 786606763\,a(n-10) + 4637371935\,a(n-11) + 5864860064\,a(n-12) - 12212688755\,a(n-13) - 42763674878\,a(n-14) - 28294485410\,a(n-15) + 57477014048\,a(n-16) + 126547355711\,a(n-17) + 64552653143\,a(n-18) - 60155915760\,a(n-19) - 96981273076\,a(n-20) - 64212382080\,a(n-21) - 123120194790\,a(n-22) - 55873185285\,a(n-23) - 22972924737\,a(n-24) + 79085877984\,a(n-25) + 44270889165\,a(n-26) + 95126815696\,a(n-27) + 10764163372\,a(n-28) + 18110185584\,a(n-29) - 9929666880\,a(n-30) - 6320643840\,a(n-31) + 1858775040\,a(n-32)
\]
for every $n>31$.
\end{theorem}

\begin{proof}
The vector $w=q(M)\tau$ was formed by $32$ matrix--vector products in exact integer
arithmetic and $\iota^{\!\top}M^{j}w$ evaluated for $j=0,1,\dots$, again exactly; those
integers vanish from the index corresponding to $n=31+1$ onwards and the run continued
until the working vector was identically zero.
\end{proof}

\section{Verification}

Four checks, each independent of the algebra above.

First, the model reproduces all $14$ terms the entry publishes, exactly, in integer
arithmetic.

Second, the count was recomputed for the smallest arrays straight from the definition:
enumerate every array of that size in the orientation the entry's own name writes and test each
cell against its neighbours directly. No window, no state and no transposition are involved,
and the published terms came back.

Third, the conjectured recurrence was evaluated directly on the published terms, with no
matrices involved, and holds wherever the proved range and the data overlap.

Fourth, the annihilation test of Lemma~\ref{lem:crit} was carried out in exact integer
arithmetic on the whole vector rather than on a sampled prefix.

\begin{thebibliography}{9}
\bibitem{oeis} The OEIS Foundation, \emph{The On-Line Encyclopedia of Integer
Sequences}, \url{https://oeis.org}, 2026. Sequence A232147.
\bibitem{stanley} R.~P.~Stanley, \emph{Enumerative Combinatorics}, Volume 1, 2nd ed.,
Cambridge University Press, 2011. (Transfer-matrix method, Section 4.7.)
\bibitem{horn} R.~A.~Horn and C.~R.~Johnson, \emph{Matrix Analysis}, 2nd ed., Cambridge
University Press, 2013.
\end{thebibliography}

\end{document}
